% !TeX root = ../../thesis.tex

\chapter*{Samenvatting}


Gebouwen kunnen gecapteerd worden met laserscanners of camera's, maar de resulterende data heeft altijd twee problemen. Ten eerste is een scan onvolledig: alles wat de scanner niet kan zien, zoals de muur achter een kast, ontbreekt. Ten tweede begrijpt een scan niet wat hij toont, want hij weet niet dat een stoel een stoel is, waardoor niets erin verplaatst, verschaald of weggehaald kan worden zoals je dat in het echt zou doen bij het herinrichten van een kamer. Hierdoor vraagt het gebruik van scans voor renovatieplanning, virtuele rondleidingen of facilitair beheer nog steeds veel manueel en duur modelleerwerk.

Dit proefschrift bouwt een automatische pijplijn die een onvolledige scan van een gemeubileerde kamer omzet in een volledige, bewerkbare 3D-omgeving, die we een \textit{Dynamic Reality Model} noemen. De pijplijn werkt stap voor stap. Eerst worden scans die op verschillende momenten of met verschillende toestellen gemaakt zijn samengevoegd tot één actuele scan, waarna bepaald wordt welke delen meubels zijn en welke delen de kamer zelf vormen, en wordt nagegaan wat écht ontbreekt versus wat er nooit was. Met die informatie worden de ontbrekende muren, vloer en plafond aangevuld met realistische geometrie en textuur, en wordt elk meubelstuk apart aangevuld, ook de delen die de scanner nooit gezien heeft. Ten slotte worden de aangevulde kamer en meubels interactief gemaakt: meubels, deuren en ramen kunnen verschaald en verplaatst worden, terwijl hun vorm realistisch blijft.

Er is geen eenvoudige manier om vanuit een echte scan na te gaan hoe een verborgen muur of stoel er écht zou moeten uitzien, dus hoorde bij dit onderzoek ook het opbouwen van een synthetische dataset van realistische virtuele kamers, V-Scan genaamd, waarin de volledige, werkelijke geometrie altijd gekend is. Dit maakte het mogelijk om de nauwkeurigheid van de methode correct te meten. De ontwikkelde pijplijn werkt betrouwbaar voor afzonderlijke, gemeubileerde kamers, en opent de deur naar snellere en goedkopere digitalisering van bestaande gebouwen voor toepassingen zoals renovaties, virtueel herinrichten, facilitair beheer, en gaming.


%%%%%%%%%%%%%%%%%%%%%%%%%%%%%%%%%%%%%%%%%%%%%%%%%%
% Keep the following \cleardoublepage at the end of this file, 
% otherwise \includeonly includes empty pages.
\cleardoublepage

\chapter*{Wetenschappelijke Samenvatting}

Digitale representaties van gebouwde binnenomgevingen vallen momenteel uiteen in twee technieken. Reality-capturemodellen, verkregen via laserscanning of fotogrammetrie, behouden geometrische en visuele nauwkeurigheid, maar blijven statische verzamelingen van punten of vlakken zonder expliciete opsplitsing in objecten. Parametrische modellen zoals Building Information Models (BIM) zijn wel bewerkbaar en semantisch gestructureerd, maar worden doorgaans op ontwerpniveau geabstraheerd en weerspiegelen de effectief gebouwde toestand niet. Geen enkele bestaande representatie combineert realisme, volledigheid en bewerkbaarheid, en er bestaat geen gevestigde pijplijn die een partiële scan rechtstreeks omzet naar zo'n representatie zonder omvangrijk manueel modelleerwerk of een terugval naar een BIM-model met beperkte visuele detaillering.

Dit proefschrift pakt die leemte aan door het Dynamic Reality Model (DRM) te introduceren: een representatie van een binnenscène als een verzameling individueel bewerkbare, objecten, ingebed in een volledige, getextureerde omgeving, met behoud van de geometrische en visuele nauwkeurigheid van de originele opname. De centrale bijdrage is een modulaire pijplijn die een partiële, heterogene scan omzet in een DRM via vijf gekoppelde stappen: uitlijning en actualisering van de scène, object- en occlusiedetectie, omgevingvervollediging, objectvervollediging en scènedynamificatie. In plaats van scène- en objectvervollediging als afzonderlijke problemen te behandelen, koppelt de pijplijn beide via expliciete occlusieredenering: de occlusiegrids die tijdens de detectiefase worden berekend, laten toe om de omgeving gericht te vervolledigen, terwijl objectvervollediging gebruikmaakt van de omliggende muren, vloer en naburige objecten als context.

Elke stap wordt kwantitatief geëvalueerd. Het uitlijnings- en actualiseringskader haalt een gemiddelde nauwkeurigheid van 0,058 m en herstelt 97,5\% van een lege scène op basis van twee partiële scans. Objectdetectie en occlusieschatting behalen een precisie van 0,495 en een recall van 0,551 bij IoU 0,5, met primaire symmetrieassen die voor zes van de zeven objectcategorieën met een fout kleiner dan 4° worden teruggevonden. Omgevingsvervollediging behaalt een Chamfer-afstand van 1,5 tot 1,9 cm, een F-score boven 0,975 bij 5 cm, en een PSNR tot 59 dB op de ingevulde texturen. Objectvervollediging behaalt, afhankelijk van de categorie, een IoU van 61,8 tot 72,2\%, waarbij een voxelgebaseerde en een beeldgebaseerde methode worden vergeleken. Omdat net de oppervlakken die het meest relevant zijn voor vervollediging ontbreken in echte scans, biedt geen bestaande dataset directe referentie voor deze evaluatie; dit proefschrift introduceert daarom V-Scan, een proceduraal gegenereerde synthetische dataset die realistische partiële scans koppelt aan volledige modellen op scène- en objectniveau.

De resultaten tonen aan dat structuur beter te herstellen is dan textuur: vlakke en geometrisch regelmatige oppervlakken worden betrouwbaar gereconstrueerd doorheen uitlijning, vervollediging en dynamificatie, terwijl textuursynthese onder sterke occlusie en onregelmatige geometrie de belangrijkste bronnen van resterende fouten blijven. De gevalideerde scope is beperkt tot eenkamerige, grotendeels vlakke, gemeubileerde binnenomgevingen; fouten door de matige detectieprecisie planten zich voort naar latere stappen, en de opsplitsing in delen voor dynamificatie blijft louter geometrisch. Binnen deze scope draagt dit proefschrift bij met een hybride, multi-bron uitlijnings- en actualiseringskader, een unificerende en modulaire scènedynamificatiepijplijn die scène- en objectvervollediging koppelt via gedeelde occlusieredenering, de V-Scan-benchmark voor gecontroleerde evaluatie van scanvervollediging, en een deelbewuste dynamificatiemethodologie die voltooide statische reconstructies omzet in interactieve DRM's.

\cleardoublepage

% vim: tw=70 nocindent expandtab foldmethod=marker foldmarker={{{}{,}{}}}
